Phase II is the confirmatory study. Its protocol was frozen and committed
before any confirmatory outcome was observed (one exception, disclosed as
deviation D1 in the protocol changelog: an 18-item pipeline smoke test whose
aggregate accuracies were observed and whose items were retained); deviations are logged in the
protocol changelog (Appendix~\ref{app:phase2-protocol}). The frozen
estimand, splits, instruments, and analysis plan are summarized here.

\subsection{Design}\label{sec:p2design}

\paragraph{Arms.} Two crossed factors: \emph{strategy forcing} (S=1: the
builder is assigned a named mechanism from a fixed set of eight; S=0: the
builder invents an improvement with no mechanism named) and the
\emph{conformance gate} (G=1: the harness must pass the frozen conformance
suite; G=0: neutral validity only). This yields four factorial arms, which
we denote \textbf{II-A} (free/ungated), \textbf{II-B} (free/gated),
\textbf{II-C} (forced/ungated), \textbf{II-D} (forced/gated). A fifth,
pre-specified \emph{secondary} arm \textbf{II-E} (open mechanism discovery:
the builder emits a self-specified behavioral contract over generic trace
primitives, gated on its own declaration) tests whether mechanism diversity
requires author-supplied vocabulary at all; it is deliberately outside the
factorial. These arm names are Phase-II-specific; Phase-I conditions in
\S\ref{sec:phase1probes} are always referred to with their descriptive
names.

\paragraph{Equal raw generation budget.} Every slot in every arm receives
exactly $R{=}3$ raw builder attempts (same model, temperature 0.7, 16{,}384
token cap). Arms differ \emph{only} in which raw candidate is admitted:
ungated arms take the first candidate passing neutral validity (parses,
imports, returns SQL on a development question --- mechanism-blind); gated
arms take the first additionally passing the conformance suite. Slot failure
is final; populations are never topped up, so generation reliability is part
of the treatment. All raw candidates are retained and generation
reliability is reported separately from population quality via three rates:
$R_{\mathrm{artifact}}$ (valid artifact), $R_{\mathrm{contract}}$ (valid,
non-vacuous contract), $R_{\mathrm{fidelity}}$ (asserted properties
verified, conditional on a valid contract).

\paragraph{Builders and seeds.} Six builders from independent labs (GLM-5.3,
Qwen3.8-Max, DeepSeek-V4-Pro, Kimi-K3, MiniMax-M3,
ERNIE-5.0-Thinking-Preview), three seeds each: 18 paired (builder $\times$
seed) cells. Every cell enters the primary contrast regardless of admitted
population size, including $K{=}0$ --- excluding low-yield cells would
condition on a post-treatment outcome.

\paragraph{Data isolation.} A fail-closed contamination registry scans every
project artifact for observed BIRD questions. Two databases touched during
development are banned wholesale; the confirmatory evaluation uses
\textbf{1{,}169 questions from nine databases never opened during development}
(the registry's defensible claim is no project-specific adaptive use, not
pretraining-time isolation; the 18-item smoke contact is the disclosed exception), with a
stratified 400-item core drawn before any outcome existed. Factorial
secondary analyses run on the same core items for all four arms.

\paragraph{Judge.} Primary: the official BIRD execution scorer (set
equality on raw result rows, 30\,s wall-clock timeout). Secondary: the
Phase-I legacy judge, measured to be uniformly conservative (on Phase-I
matrices it never credits an answer the official scorer rejects).

\paragraph{Instruments.} The conformance suite drives each harness against
scripted scenarios and asserts on the recorded execution trace, never on
task performance, so gating cannot leak outcome supervision. Every mechanism
is probed with \emph{paired counterfactual variants} whose correct answer
differs --- including a variant where the first query \emph{succeeds},
separating conditional repair from an unconditional second call --- and
answers are compared by exact normalized equality. The suite was calibrated
before any Phase-II generation against three hand-written positive controls
(all must pass) and two known-negative impostor harnesses (both must be
rejected), plus the invariant that the bare baseline must violate every
conditional property; the calibration is an executable gate re-run before
each generation batch. Arm II-E additionally enforces non-vacuity (the
contract must be violated by the bare baseline, checked by running it) and
verifies only the properties actually asserted.

\begin{figure}[t]
\centering
\includegraphics[width=0.98\linewidth]{figures/fig_design.pdf}
\caption{\textbf{Phase-II confirmatory design} (frozen protocol). Left:
database-level contamination isolation (the two development databases are
banned wholesale; evaluation uses nine databases never opened during development). Middle:
the $2{\times}2$ factorial (strategy forcing $\times$ conformance gate)
with equal raw budgets, plus the open-mechanism arm II-E. Right: the
three-layer admission pipeline (neutral validity $\rightarrow$
counterfactual conformance $\rightarrow$ non-vacuity vs.\ bare).}
\label{fig:design}
\end{figure}

\subsection{Frozen analysis plan}\label{sec:p2plan}
The \textbf{primary estimand} is the bare-inclusive oracle headroom $H(P)$
(Equation in \S\ref{sec:metrics}; $K{=}0$ gives $H=0$, and $K{=}1$ is a
genuine bare-vs-harness routing problem) of arm II-D minus arm II-A,
paired by (builder $\times$ seed), on the full 1{,}169-question set under
the official judge --- a continuous effect with two-sided 95\% CI and a
paired permutation test, not a pass/fail threshold. The Phase-I thresholds
(union repair $\geq$15\%, headroom $\geq$5\,pp) are retained only as a
descriptive operational checklist; threshold-crossing itself is
statistically uncertain and is never used as a scientific verdict.

\textbf{Confirmatory secondaries} (Holm-corrected, tested after the
primary, all four arms on the same core-400 items): gate main effect
$[(B{-}A)+(D{-}C)]/2$; strategy main effect $[(C{-}A)+(D{-}B)]/2$; and the
strategy $\times$ gate interaction $(D{-}C)-(B{-}A)$ --- the key
mechanism-level question of whether strategy differentiation requires
execution fidelity to be realized.

\textbf{Inference} uses a hierarchical paired bootstrap (10{,}000
replicates): databases resampled with replacement, then tasks within each
drawn database, with (builder $\times$ seed) cells resampled as paired
units. K-matched sensitivity ($k{=}1\ldots K_{\min}$) and a
population-conditional analysis ($K \geq 3$ both arms) are reported as
sensitivity checks. Note that $H(P)$ is \emph{not} monotone in $K$: oracle
accuracy is non-decreasing in $K$ by construction, but so may be the best
fixed accuracy --- adding a harness that answers everything correctly
drives headroom to zero --- which is why matched-K comparison, not
monotonicity, licenses cross-population claims.

\textbf{Robustness audits (frozen):} dual-judge scoring on every cell;
cache-off replication of arms A and D on the core (quantifying how much
apparent agreement is the shared response cache, i.e.\ common random
numbers, versus structural); and three independent cache-off repeats of
bare and arm D on the core, giving an outcome flip rate and repair
stability rate.
